\section{A complementary degree-three source and the exact free-pair test}
\label{sec:136-degree-three-common-skeleton}

\paragraph{Archive scope.}
The vertical and horizontal conditions are cote136, physical
pp.~117 and 119, footers116 and118, scans117 and119. The
common-source question is physical p.~121, footer120, scan121.
The construction below complements ARCH-136-DEGREE-TWO-FREE:
it preserves a common first skeleton and exact cohomology
through degree two, while concentrating the remaining integral
vertical problem in degree three. It does not make the main
common-source assertion stronger than the already established
reductions. The exact degree-three attachment criterion explains
one of its remaining restrictions.

\paragraph{ARCH-136-Q3-COEFFICIENT-KERNEL.}
Let \(B\) be a nonnegative strict PD differential algebra over
\(S=\Gamma_{\mathbb Z}(T)\), with \(B^0\) torsion-free as
an abelian group, and suppose that in every weight
\[
             H^0(B)=S,\qquad H^1(B)=H^2(B)=0.             \tag{136.QK1}
\]
The degree-zero identification is the specified coefficient map.
Choose any family of homogeneous elements and integers
\[
 z_i\in Z^3(B)^{m_i},\quad b_i\in B^{m_i,2},\quad
 N_i>0,\qquad db_i=N_i z_i.
\]
Form the free strict PD extension with even \(c_i\) in form
degree two and odd \(t_i\) in form degree one, both of weight
\(m_i\), and with
\[
                  dc_i=z_i,\qquad dt_i=N_i c_i-b_i.
                                                               \tag{136.QK2}
\]
Then \(H^0(E)=S\), \(H^1(E)=0\), and there is a canonical
graded \(S\)-module identification
\[
 H^2(E)\ \cong\
 \ker\left(\bigoplus_i(S/N_iS)\langle m_i\rangle
                \longrightarrow H^3(B),\quad
                (\overline s_i)_i\longmapsto
                            \sum_i s_i[z_i]\right).
                                                               \tag{136.QK3}
\]
Here the weight-shift notation places the distinguished generator
in weight \(m_i\). The sum has finite support. Thus such a
free attachment preserves \(H^2=0\) exactly when its selected
cyclic coefficient modules inject into \(H^3(B)\).

\begin{proof}
The monomial basis of the free extension makes the groups in
degrees at most two explicit. A degree-one cycle is
\(x+\sum_i f_i t_i\), with \(x\in B^1\), \(f_i\in B^0\).
The coefficient of the independent new vector \(c_i\) in its
differential is \(N_i f_i\). Torsion-freeness of \(B^0\)
forces every \(f_i=0\). Its cycle and boundary groups are
therefore the old ones. Form degree zero is unchanged, proving
the first two assertions.

A general degree-two element has a unique expression
\[
 x+\sum_i y_i t_i+\sum_i w_i c_i
                      +\sum_{i<j}k_{ij}t_it_j,
 \quad x\in B^2,\ y_i\in B^1,\ w_i,k_{ij}\in B^0.
                                                               \tag{136.QK4}
\]
For a cycle, the independent \(c_i t_j\) coefficients of
its differential force \(N_i k_{ij}=0\), hence
\(k_{ij}=0\). There are no other terms of the differential
having those monomials. The remaining cycle equations are
\[
 dy_i=0,\qquad dw_i=N_i y_i,\qquad
             dx+\sum_i y_i b_i+\sum_i w_i z_i=0.          \tag{136.QK5}
\]
Use \(H^1(B)=0\) to write \(y_i=dv_i\), with
\(v_i\in B^0\). The second equation and \(H^0(B)=S\)
then give \(w_i=N_i v_i+s_i\), for a coefficient
\(s_i\in S\) in the appropriate weight. Subtracting the
boundary \(\sum_i d(v_i t_i)\) reduces (136.QK4) to
\[
                     x'+\sum_i s_i c_i,
                 \qquad dx'=-\sum_i s_i z_i.             \tag{136.QK6}
\]
Consequently \((\overline s_i)_i\) belongs to the kernel
in (136.QK3). Changing a chosen primitive \(v_i\) changes
\(s_i\) by a multiple of \(N_i\), so the class is
well-defined. Adding the differential of a general degree-one
element makes the same change; its old part contributes only
an old boundary.

Conversely every tuple in that kernel admits an \(x'\)
satisfying (136.QK6), by the definition of \(H^3(B)\), and
hence comes from a cycle. If \(s_i=N_i r_i\) for every
\(i\), subtract \(\sum_i d(r_i t_i)\). The remainder is
an old closed element of \(B^2\), which is an old boundary
because \(H^2(B)=0\). This proves injectivity as well as
surjectivity. All changes are homogeneous, coefficient-linear,
and finite on each element, so the identification has exactly
the stated graded and arbitrary-family scope.
\end{proof}

\paragraph{What the criterion does and does not remove.}
For one native class generating an embedded copy of
\(S/NS\), (136.QK3) is zero. A coefficient-tail multiple of
such a class need not have this property: over \(S/pS\)
every positive-weight element is nilpotent, so multiplying the
class by a nonzero positive-weight coefficient produces a
nontrivial annihilator. Formula (136.QK3) then creates new
\(H^2\), even though the input class is nonzero and has its
correct integer order. This is the degree-three version of
ARCH-136-TAIL-MOORE, with no redundant-pair assumption.

This is a pointwise cohomology test, separate from the obstruction
to selecting primitives with common first-skeleton restrictions.
It cannot remove the latter by itself. For example, the
already checked ARCH-136-SKELETON-OBSTRUCTION with \(d=2\)
has native coefficient-independent degree-three classes
\((S/pS)\otimes U_n\), but its two prescribed maps agree on
the entire first skeleton and still have a nonzero two-sphere
residue. That example is recalled only to delimit the criterion;
it is not an obstruction to freely choosing the desired common
source and is not a new counterexample here.

\paragraph{ARCH-136-DEGREE-THREE-COMMON-SKELETON.}
For every finite nonempty family of integral admissible theories,
there are their admissible replacements \(C'_j\to C_j\) from
AQ-136-FIRST-SKELETON and one strict source \(A\to C'_j\)
with the following properties:
\[
\begin{gathered}
 \operatorname{sk}_1 A=F=\operatorname{sk}_1 C'_j,
 \qquad A\to C'_j\text{ is the identity on }F,\\
 A_n^{w,q}\text{ is free abelian for every }n,w,q,\\
 H^0(A_n^{w,*})=S^w,\qquad
                 H^1(A_n^{w,*})=H^2(A_n^{w,*})=0,\\
 H^*((A_n^{w,*})_{\mathbb Q})=S_{\mathbb Q}^w,\\
 H_r(NA_*^{w,q})=0\quad(w\geq\mu_q,\ r\geq0),\\
 H^q(A_n^{w,*})=0\quad(q\geq4,\ w\geq\lambda_q,\ n\geq0).
\end{gathered}                                             \tag{136.DT3C1}
\]
The bounds \(\mu_q,\lambda_q\) are fixed during the construction.
The only missing vertical condition is eventual vanishing of
the integer-torsion group \(H^3\). No uniform finite-prime
support or prescribed-skeleton completion theorem is asserted.

\begin{proof}
Take the common first skeleton \(F\) from
AQ-136-FIRST-SKELETON, with its generators \(x_m,e_m\),
\(m\geq N\), in forms zero and one and \(dx_m=e_m\).
Start from its simplicial left Kan extension from degrees
zero and one. In degree \(n\) its nonconstant generators
are the copies of these free even unit disks indexed by the
surjections \([n]\twoheadrightarrow[1]\); constant maps
give the coefficient copies, which are identified. Thus it
is, pointwise, a free PD extension of \(S\) by even unit
disks. This also follows directly by presenting the left Kan
extension as the coproduct of the copies of \(F_1\) over
their common coefficient algebra. The construction has exactly
the specified first skeleton, is free abelian weightwise,
and has coefficient vertical cohomology integrally in every
simplex. The left Kan universal property supplies its maps
to every \(C'_j\), extending the identity on \(F\).

There is no horizontal degree-zero obstruction in the rows
chosen for repair. In positive form degrees the simplicial
degree-zero module is zero. In form degree zero and weight
\(w\geq N\), the generator \(x_w\in F_1\) already has
boundary \(T^{[w]}\). Hence every remaining horizontal
cycle can be filled by free horizontal cells of simplex
dimension at least two. Choose the row cutoffs once, above
the target cutoffs and with \(\mu_0\geq N\). Attach the
same free vertical unit-disk cells as in
ARCH-136-DEGREE-TWO-FREE, now only in these dimensions.
The target row extension argument supplies all their images;
these attachments leave the entire first skeleton unchanged.

Alternate this horizontal step with the free Tate-pair step
only in form degrees \(q\geq4\). Rational vertical exactness
ensures that an old closed \(z\) of form degree \(q\) has
an integer relation \(db=N'z\), with \(N'>0\). In a weight
above the fixed target bounds, attach generators \(c,t\)
of forms \(q-1,q-2\) with \(dc=z\), \(dt=N'c-b\).
Use ARCH-136-SKELETON-TATE-LIFTS in its \(q\geq4\) case:
one can choose a common solution on
\(G_n=\operatorname{sk}_1\Delta[n]\), and extend it to
solutions in all targets. Thus the attachment is made relative
to that graph, with the common boundary values in \(F\).
Here the source itself has first skeleton exactly \(F\), so
these common graph values are existing source elements.
Identifying the graph generators with those elements gives
a differential-stable PD ideal, since the common values obey
the same two differential equations. This permits the full
graph identification, rather than only a vertex identification.
Its new generator modules in each simplex are free on the
relative simplexes \(\Delta[n]/G_n\), and its maps to every
target continue to be the identity on \(F\).

The bounds for this last use are uniform in \(n\): the graph
cone calculation in ARCH-136-SKELETON-TATE-LIFTS has no
cohomology in degrees at least two, and the required extension
uses only finitely many fixed target row and vertical cutoffs
for each form degree \(q\). Denote their maximum by
\(\lambda_q\), and retain it at every stage.

We check the low-degree assertion rather than importing it
from the degree-two construction. A free horizontal even disk
is integrally contractible. After the arity-one contractible
block of a free odd unit disk is removed, every remaining
positive-arity block starts in form degree at least three. The
coefficient contractions and the nonnegative grading therefore
show that horizontal attachments preserve \(H^0,H^1,H^2\).
For a Tate pair with \(q\geq5\), no new generator has form
degree below three, so these same groups are unchanged.
For \(q=4\), write the new independent generators as
\(t_\alpha\) of form two and \(c_\alpha\) of form three.
A new degree-two cycle is
\(x+\sum_\alpha f_\alpha t_\alpha\), with
\(f_\alpha\in A^0\). Its differential has coefficient
\(N'_\alpha f_\alpha\) on the independent \(c_\alpha\).
The old degree-zero module is torsion-free, so all
\(f_\alpha=0\). Its cycles are old and its boundaries
are old. This proves preservation of \(H^2=0\), while
degrees zero and one are visibly unchanged. The argument is
unaffected by prescribed graph boundary values: after choosing
the free relative basis they change only old terms in the
differential, not the leading \(N'_\alpha c_\alpha\).

Every step remains free abelian in each bidegree. Pointwise
it is a free graded PD extension on a degreewise free relative
parameter module; its old graded inclusion is split by setting
the new independent generators to zero. Compatible monomial
bases give freeness at union stages as well. In particular the
torsion-freeness used above is available throughout.
Each horizontal disk and each Tate pair is rationally
contractible relative to the old algebra. For the latter the
change \(c'=c-b/N'\) gives \(dc'=0\), \(dt=N'c'\).
Hence rational vertical coefficient exactness holds at all
stages, and the remaining integral positive cohomology is
integer torsion.

Repeat the two steps countably many times, attaching all
indicated cycles at each step. Every element has finite
support. A horizontal cycle in a chosen row tail is therefore
killed at a later horizontal step; a vertical class in degree
\(q\geq4\) and weight \(w\geq\lambda_q\) is killed at
a later vertical step. Filtered colimits are exact and commute
with the kernel/cokernel computations in fixed degrees. Thus
the same fixed bounds give precisely (136.DT3C1) in the
colimit. No degree-three Tate step was used. This proves
all the assertions and the stated limitation.
\end{proof}

\paragraph{Remaining comparison problem.}
For a further degree-three pair, preserving the common first
skeleton requires the graph class in
\(H^1(G_n;S^w/N S^w)\) from
ARCH-136-SKELETON-TATE-LIFTS to vanish. Even after such a
choice is possible, the new low cohomology must meet the
coefficient-module injectivity test (136.QK3) if one wishes
to retain \(H^2=0\). The present construction proves neither
condition automatically. It is an alternative location of the
remaining compatibility problem, not a claimed resolution of it.
